% ------------------------------------------------------------------
%  Discussion (~900 words). Interpretation, why structure beats prompting,
%  limitations stated plainly, and what a prospective study must show.
% ------------------------------------------------------------------

\mngs has made the pathogen list longer, not the diagnosis clearer. The results here show that the interpretive step can be made explicit, reproducible and auditable without giving up the sensitivity that justifies ordering the test. On \todo{55} patients from three hospitals, a pipeline of modality-specific extraction, deterministic scoring, language-model safety review and case-matched literature rescue recovered \todo{91\%} of physician-adjudicated pathogens at a precision of \todo{72\%}, whereas the raw report, the multiplex panel and culture each delivered one of those two properties at the expense of the other. The same pipeline exceeded four frontier language models given identical data by \todo{14--15} F1 points.

Why does structure beat prompting? The direct-prompting baselines were not naive: the prompt stated the colonization rules, the site rules and the meaning of the control codes, and the models reasoned at medium effort with a strict output schema. Their precision matched the pipeline's. What they lost was recall, and the pattern of their misses is informative. Asked to be cautious about \textit{Candida}, herpesviruses and environmental organisms, a model becomes cautious about everything, and declines the \textit{Stenotrophomonas} with dominant reads in a ventilated patient or the \textit{Pneumocystis} with modest reads in a patient on corticosteroids that the rules, the host tier and the literature together support \todo{replace with real examples from the per-patient comparison}. A rule can encode the exception precisely---R2 reads suffice for a protected pathogen at rank~1, but not for a water-borne Gram-negative without dominance---whereas a prompt can only ask for judgement. The literature stage adds what neither a rule nor a general prompt has: the actual evidence that this organism causes this syndrome at this site in this kind of host, judged article by article with the span that supports it. That the newest model (\gptastra) improved on its predecessors by only \todo{2.6} recall points suggests the ceiling is set by the task formulation, not by model capacity.

The design choices that follow from auditability are worth stating as principles, because they apply beyond \mngs. First, language models extract and review; rules decide. Every organism that reaches the physician passed a rule whose inputs are recorded, so the decision can be replayed. Second, the answer key is quarantined: reference labels enter only the evaluation code, the case-fit prompt strips upstream picks, reads, cultures and review rationales, and the direct-prompting payloads were audited for label leakage before any model was called. Third, selection is frozen before explanation. Rationales are generated afterwards, restricted to the evidence window and checked for structure, so an eloquent paragraph cannot promote an organism the rules rejected. Fourth, upstream picks are immutable and rescue only appends. This bounds the precision of the final list by that of the deterministic stage and makes the contribution of each downstream module measurable: here, rescue changed \todo{7\%} of selections. These properties cost something---the pipeline calls a language model at four points and PubMed at one, and a full run on a new patient takes minutes and money rather than seconds---but they are what turns a research prototype into something a laboratory could validate.

Several limitations bound the claims. The study is retrospective and the cohort is small: \todo{55} patients for the main comparison and \todo{41} for the model comparison, from three hospitals in one country, sequenced by \todo{one} laboratory. The reference standard is physician adjudication informed by all available results, including \mngs, so it is not independent of the test being evaluated, and its labels were revised in versioned steps during the study; we report the version used and the revision history (Methods), but a fully blinded reference would be stronger. The evaluation cohorts contain only patients with at least one adjudicated pathogen, so our precision estimates do not cover the false-positive risk in patients with no infection, which is where a pipeline built to preserve recall is most likely to err. The scoring rules and guardrails were developed by iterating on these same cases with the labels visible; the \todo{55}-patient result is therefore a development-set result, not an untouched external test, and the gain over direct prompting should be read as a comparison of formulations on a shared cohort rather than as an estimate of prospective performance. The two vignettes and the organism-burden comparison come from an earlier \todo{35}-patient series evaluated with the first version of the pipeline \todo{confirm}. The literature and case-fit stages depend on a commercial language model and on live PubMed; frozen article snapshots make the reported decisions replayable, but a rerun on a new patient will retrieve different articles as the literature grows, and the case-fit judgements, although constrained by a traceability gate, are not deterministic. Finally, we measured agreement with adjudication, not clinical benefit: whether presenting two organisms instead of six changes antimicrobial choice, time to appropriate therapy or outcome remains to be shown.

These limitations define the next study. A prospective evaluation should freeze the rules and prompts before enrolment, include patients without infection, use an adjudication committee blinded to the pipeline output, and measure antimicrobial changes as an endpoint. Several extensions are straightforward within the present architecture. The reference-label revisions and per-patient error analysis point to rule refinements that can be versioned and re-evaluated without touching the orchestrator. Antimicrobial-resistance genes on the same \mngs report, already parsed by the front end, could enter as a further evidence module. And because the \mngs signal features, the site logic and the literature stage are organism-agnostic, the pipeline can be extended to bloodstream and central nervous system infection, where the colonization problem is smaller but the stakes of a missed pathogen are higher. \mngs was meant to end the era of pneumonia treated without a diagnosis; a pipeline that can say which of the organisms it found is the diagnosis, and show why, brings that closer.
